% Generated draft: requires mathematical and source review.
\documentclass[11pt]{article}
\usepackage[margin=25mm]{geometry}
\usepackage{fontspec}
\setmainfont{Latin Modern Roman}
\usepackage{amsmath,amssymb,mathtools,unicode-math}
\setmathfont{Latin Modern Math}
\usepackage{enumitem,needspace,xurl,hyperref}
\setlength{\parindent}{0pt}
\setlength{\parskip}{5pt}
\setlength{\emergencystretch}{3em}
\newcommand{\N}{\mathbb{N}}
\newcommand{\Z}{\mathbb{Z}}
\newcommand{\Q}{\mathbb{Q}}
\newcommand{\R}{\mathbb{R}}
\newcommand{\C}{\mathbb{C}}
\newcommand{\F}{\mathbb{F}}
\newcommand{\A}{\mathbb{A}}
\newcommand{\PP}{\mathbb{P}}
\newcommand{\E}{\mathbb{E}}
\newcommand{\eps}{\varepsilon}
\newcommand{\dd}{\,\mathrm{d}}
\begin{document}
\section*{UnsolvedMath ID: 30001697}
\textbf{Source problem code:} catalogue entry
\subsection{Definitions and mathematical statement}
For a positive integer \(n\), let
\[
S^{n-1}=\{x\in\R^n:\|x\|_2=1\},\qquad
O(n)=\{\rho\in\R^{n\times n}:\rho^{\mathsf T}\rho=I_n\}.
\]
Here \(O(n)\) includes both orientation-preserving and orientation-reversing orthogonal transformations; the source uses the word ``rotation'' for an element of this group.

Interpreting the point count \(\ell\) as a positive integer, the supplied question asks whether
\[
\begin{gathered}
\forall\ell\in\Z_{\geq1}\ \exists n\in\Z_{\geq1}\ 
\forall x_1,\ldots,x_\ell\in S^{n-1}\text{ pairwise distinct}\ 
\forall f\in C(S^{n-1},\R),\\
\exists\rho\in O(n)\quad
f(\rho x_1)=\cdots=f(\rho x_\ell).
\end{gathered}
\]
The dimension \(n\) may depend on \(\ell\), but must be chosen independently of the point set and the continuous function. The transformation \(\rho\) may depend on both.

\textbf{Unresolved qualification.} The source supplies no lower bound on \(n\) beyond that needed to interpret the sphere. Since \(S^0\) has only two points, choosing \(n=1\) makes the universal assertion over \(\ell\)-element subsets empty for \(\ell\geq3\). Thus the displayed literal formulation has a vacuity that conflicts with the nontrivial general finiteness question described in the supplied background. A restriction ensuring that such point sets exist, or another intended dimension restriction, requires source confirmation; none is silently added here.
\subsection{Short English statement}
For each positive point count \(\ell\), does a dimension exist in which every \(\ell\)-point subset of the unit sphere can be moved by an orthogonal transformation into one level set of every continuous real-valued function on that sphere? The supplied formulation needs clarification of its dimension restriction to exclude vacuous choices.
\subsection{Sources}
\begin{itemize}
\item Local UnsolvedMath registry: supplied record from \texttt{data/unsolvedmath/problems.json}, permanent ID \(30001697\), title ``Weak Knaster Conjecture,'' exact source problem code \texttt{catalogue entry}. The record attributes the problem to the 2011 Oberwolfach workshop collection and identifies no individual proposer.
\item Supplied source citation: \emph{Topological and Geometric Combinatorics} (2011), Oberwolfach Reports, DOI 10.4171/owr/2011/08, pp.~349--423. Original OWR report: \url{https://doi.org/10.4171/owr/2011/08}.
\item Supplied literature reference, labeled ``Bukh--Karasev weak Knaster results'': \url{https://arxiv.org/abs/1212.5351}.
\item Literature assessment checked on \texttt{2026-08-21}: the supplied curation labels the problem partially solved, reports substantial special-case results, and states that the general finiteness assertion remains open and that \(n(4)\) was unresolved in the primary literature located. This is supplied dated curation derived from the OWR collection, not an independently verified literature assessment. Its time-sensitive status and cited claims require independent verification.
\item The registry labels the statement \texttt{corrected\_verified} and reports a manual semantic, typographic, and LaTeX review. These are supplied editorial claims, not independent mathematical verification of this draft.
\end{itemize}
\end{document}
